\lecture{14}{23. Oktober 2025}{Torsional Deformation}


\section{Torsion}

\subsection{Torsional Deformation of a Circular Shaft}
\begin{definition}[Torque]
  Torque is a moment that tends to twist a member about its longtitudinal axis.
\end{definition}


We consider a torque applied to a circular shaft made of a highly deformable material such as rubber.
As torque is applied, the longtitudinal grid lines tend to distort into a helix as shown on \Cref{fig:f14_1}.
Also, all the cross sections of the shaft will remain flat and all radial lines will remain straight and simply rotate during the deformation.
For a small angle of twist the length of the shaft and its radius will remain unchanged. 

\begin{figure} [ht]
  \centering
  \includegraphics[width=0.85\linewidth]{./figures/f14_1.png}
  \caption{Distortion of grid lines in a highly deformable material due to an applied torque.}
  \label{fig:f14_1}
\end{figure}

If we now fix one end of the shaft and apply a torque to the other end, then the undeformmed plane in \Cref{fig:f14_2}.a will distort into a skewed plane as shown.

A radial line, located on the cross section at distance $x$ from the fixed end of the shaft, will rotate through an angle $\phi(x)$.
This angle is called the \textit{angle of twist}.
As can be seen from the arguments to $\phi$ its magnitude depends on the position $x$ along the shaft.

\begin{figure} [ht]
  \centering
  \includegraphics[width=0.85\linewidth]{./figures/f14_2.png}
  \caption{Torque applied to a shaft fixed at the other end.}
  \label{fig:f14_2}
\end{figure}


We will now consider a small disk-element, located at $x$ as shown on \Cref{fig:f14_2}.b.
Due to the deformation, the both the front and rear faces of the element will undergo rotation -- the back face by $\phi(x)$ and the front face by an infinitesimally larger magnitude $\phi(x) + \mathrm{d}\phi$.
The difference between these two rotations $\mathrm{d}\phi$ between the front and back face of the element causes a shear strain $\gamma$.

This angle (or shear strain) can be related to the angle $\mathrm{d}\phi$ by noting that for small angles $s = \theta r$, hence the length of the red arc in \Cref{fig:f14_2}.b is
\begin{equation} \label{eq:101}
  \rho \, \mathrm{d}\phi = \mathrm{d}x \gamma \implies \gamma = \rho \frac{\mathrm{d}\phi}{\mathrm{d}x} 
\end{equation}
Since $\gamma_{\mathrm{max}}$ occurs for $\rho = c$, then $\mathrm{d}\phi / \mathrm{d}x = \gamma / \rho = \gamma_{\mathrm{max}} / c$ and therefore:
\[ 
\gamma = \frac{\rho}{c} \gamma_{\mathrm{max}}
.\]
I.e. the shear strain varies linearly along the radius of a cross section from zero at the center to a maximum $\gamma_{\mathrm{max}}$ at the outer edge.


\subsection{The Torsion Formula}
When an external torque is applied to a shaft it creates a corresponding internal torque $T$ within the shaft.
Here an equation that relates this internal torque to the shear stress distribution acting on the cross section of the shaft will be found.

If the material is linear elastic, then Hooke's law applies, i.e. $\tau = G \gamma$, or $\tau_{\mathrm{max}} = G \gamma_{\mathrm{max}}$.
Consequently, a linear variation in shear strain leads to a linear variation in shear stress along the radius of a cross section.
Hence, $\tau$ will vary from zero at the shafts longtitudinal axis to a maximum $\tau_{\mathrm{max}}$ at its outer edge.
\begin{equation} \label{eq:103}
  \tau = \frac{\rho}{c} \tau_{\mathrm{max}}
\end{equation}

\begin{figure} [ht]
  \centering
  \includegraphics[width=0.35\linewidth]{./figures/f14_3.png}
  \caption{Linear variation of shear stress along a radial line.}
  \label{fig:f14_3}
\end{figure}

On \Cref{fig:f14_3}, each element of the area, $\mathrm{d}A$, located at $\rho$ is subjected to a force of
\[ 
\mathrm{d}F = \tau \, \mathrm{d}A
.\]
The torque produced by this force is
\[ 
\mathrm{d}T = \rho \left( \tau \, \mathrm{d}A \right)
.\]
For the entire cross section we require
\[ 
T = \int_{\mathrm{A}} \rho \left( \tau \, \mathrm{d}A \right) = \int_{\mathrm{A}} \rho \frac{\rho}{c} \tau_{\mathrm{max}} \, \mathrm{d}A
.\]
Since $\tau_{\mathrm{max}} / c$ is constant we can rewrite this as:
\[ 
T = \frac{\tau_{\mathrm{max}}}{c} \int_{\mathrm{A}} \rho^2 \, \mathrm{d}A
.\]

This integral represents the polar moment of inertia of the circular cross section about the shaft's longtitudinal axis.
This is symbolized as $J$.
As a result the above equation can be rewritten as:
\begin{equation} \label{eq:106}
  \tau_{\mathrm{max}} = \frac{Tc}{J}
\end{equation}
Where:
\begin{center}
   \begin{tabular}{ll}
      $\tau_{\mathrm{max}}$  & The maximum shear stress within the shaft, which occurs at its outer surface \\
      $T$ & The resultant \textit{internal torque} acting at the cross section. \\
      $J$ & The polar moment of inertia of the cross-sectional area \\
      $c$ & The outer radius of the shaft
  .\end{tabular}
\end{center}

We can insert \Cref{eq:106} into \Cref{eq:103} to get the shear stress at the intermediate distance $\rho$ on the cross section as
\[ 
\tau = \frac{T \rho}{J}
.\]

Either of the above two equation are referred to as the torsion formula.
This is only applicable if the shaft has a circular cross section, the material is homogeneous, and behaves in a linear elastic manner.


\subsection{Power Transmission}
Shafts and tubes with circular cross sections are often used to transmit power from a motor.
The work transmitted by a rotating shaft is equal to the torque applied times the angle of rotation.
Therefore, if during an instant of time $\mathrm{d}t$ an applied torque $\Vec{T}c$ causes the shaft to rotate $\mathrm{d}\theta$, then the work done is $T \, \mathrm{d}\theta$ and the instantaneous power is
\[ 
P = T \frac{\mathrm{d}\theta}{\mathrm{d}t} 
.\]
The shaft's angular velocity is $\omega = \mathrm{d}\theta / \mathrm{d}t$, and hence the power becomes:
\[ 
P = T \omega
.\]

For motors the frequency of a shaft's rotation $f$ is often reported.
Using this the power is
\[ 
P = 2 \pi f T
.\]
